\documentclass[../../main.tex]{subfiles}

\begin{document}

\section{Infinite Descent}

Infinite descent is a powerful proving tool in number theory that
is used very frequently. Many, if not most, problems on Diophantine
equations are proven with infinite descent and it has wide range
of applications from showing that there are no solutions under
a certain condition and finding solutions to advance techniques
like Vieta jumping.

Before we get started with how this technique works, we need to
discuss the foundational principle that infinite descent builds
up on. Consider the following principle.

\begin{theorem}{}{}
The \textbf{well-ordering principle} states that there always
exists a least element in an non-empty set of positive integers.
\end{theorem}

In olympiad, we take the principle above as an axiom and does not
require a proof. Meaning, we can just use ``by well-ordering
principle...'' when we prove by infinite descent.

The well-ordering principle is foundational for infinite descent
because it allows us to have contradiction. Here is what I mean.
When we prove with infinite descent, we assume that there exists
a minimum value in a set. The minimum value could be a number,
sum, and much more depending on the problem. Also, note that we
can assume this because our interests only lies on positive
integers and there cannot be a ``least element'' for rational
and reals. Then, we will infinitely descend, eventually getting
contradiction with our assumption. This is quite hard to understand
without an example, so let's take a look at very famous example
before formally stating the method described.

\begin{problem}{Classic Problem}{}
Show that the equation $x^3 + 2y^3 = 4z^3$ does not have any
solutions in positive integers.
\end{problem}
(\href{https://youtu.be/UJR56wT0NIM}{Video Solution})
\begin{solution}
For the sake of contradiction, assume that there exists positive
integer solutions to the equation and let $(a, b, c)$ be the
solution with the least $a$ by well-ordering principle.
Consider the following congruence.
\[
x^3 + 2y^3 \equiv 4z^3 \equiv 0 \pmod{2}
\]
Therefore, $2 \mid x$. Let $x'$ be an integer such that $x = 2x'$.
Substituting, $8x'^3 + 2y^3 = 4z^3$ and $4x'^3 + y^3 = 2z^3$ are
obtained. By the similar argument, it is evident that $2 \mid y$
and $y = 2y'$ for some integer $y'$. Substituting again,
$4x'^3 + 8y'^3 = 2z^3$ and $2x'^3 + 4y'^3 = z^3$. Therefore,
$2 \mid z$ and there exists some positive integer $z'$ such that
$z = 2z'$. Continuing with the final substitution,
$2x'^3 + 4y'^3 = 8z'^3$ and $x'^3 + 2y'^3 = 4z'^3$ are obtained.
In other words, if $(x, y, z)$ is a solution, $(x', y', z')$ is
also a solution to the equation.

Note that $(a, b, c)$ is assumed to be a solution with the least $a$.
Therefore, $(\frac{a}{2}, \frac{b}{2}, \frac{c}{2})$ is also a
solution to the equation. Because $\frac{a}{2} < a$, there is no
positive integer solutions to the equation.
\end{solution}

As you can see from the example above, we assume the minimum
and show that the infinite descent contradicts with our assumption.

\begin{definition}{}{}
The \textbf{proof by infinite descent} is a technique that shows
contradiction through infinitely descending sequence. Assume that
there exists a solution with minimum condition by well-ordering
principle. Then, construct a strictly and infinitely descending
sequence. The proof completes by contradiction.
\end{definition}








\end{document}


